\textbf{结论名称：}含参不等式恒成立——分类讨论法（参数不可分离时的变号因子处理）

\textbf{适用条件：}形如 \(a\cdot k(x) \ge h(x)\)（或 \(\le\)）在区间 \(D\) 上恒成立，且 \(k(x)\) 在 \(D\) 上变号，无法统一分离参数。

\textbf{变量说明：} \(a\) 为待求参数；\(k(x),\,h(x)\) 为定义在 \(D\) 上的已知函数；按 \(k(x)\) 的符号将 \(D\) 划分为 \(D_1=\{x\in D\mid k(x)>0\}\)，\(D_2=\{x\in D\mid k(x)<0\}\)，\(D_0=\{x\in D\mid k(x)=0\}\)。

\textbf{核心公式：}
对于 \(\forall x\in D,\, a\cdot k(x)\ge h(x)\)，等价于分类讨论：
\[
\boxed{
\begin{cases}
a\ge \dfrac{h(x)}{k(x)}, & x\in D_1 \\[6pt]
a\le \dfrac{h(x)}{k(x)}, & x\in D_2 \\[6pt]
0\ge h(x), & x\in D_0
\end{cases}
}
\]
分别求出使各部分恒成立的 \(a\) 的范围，最后取交集即得整体恒成立的 \(a\)。

\textbf{使用场景：}导数综合题中，含参不等式无法直接分离参数（参数系数变号）；或分离后函数形式复杂，命题期望用分类讨论求解。多见于含 \(a(x-c)\)、\(a\ln x\) 等因子在区间上跨零点的问题。

\textbf{简单例子：}已知不等式 \(a(x-1)\ge x^2-1\) 对一切 \(x\in[0,2]\) 恒成立，求 \(a\) 的取值范围。

解：变形为 \(a(x-1)\ge (x-1)(x+1)\)，按 \(x-1\) 的符号讨论：\\
\(x=1\)时，\(0\ge 0\)，恒成立；\\
\(x\in[0,1)\)时，\(x-1<0\)，得 \(a\le x+1\)，需 \(a\le \min_{[0,1)}(x+1)=1\)；\\
\(x\in(1,2]\)时，\(x-1>0\)，得 \(a\ge x+1\)，需 \(a\ge \max_{(1,2]}(x+1)=3\)。\\
同时满足 \(a\le 1\) 且 \(a\ge 3\)，无解，故不存在满足条件的 \(a\)。